\documentclass[11pt]{article}

\usepackage[margin=1in]{geometry}
\usepackage{amsmath, amssymb}
\usepackage{booktabs}
\usepackage{hyperref}

\title{Endpoint construction}
\author{Tijn Jacobs}
\date{\today}

\begin{document}
\maketitle

\subsection{Endpoint construction}

The endpoint is deliberately asymmetric across the two data sources. The
randomized component (ACTG~175) is used as published: a single composite
indicator and one follow-up time, with the underlying components not
separately released. The observational component (MACS) is reconstructed
component-wise from the public release and then assembled to match a
chosen target endpoint. Five candidate assemblies are tabulated below.
The Kaplan--Meier diagnostic figure uses variant v5 (overall survival)
for signal quality, while the publication-target FusionForest fit uses
variant v3 (\emph{EFS-cd4-fix}), the closest analogue of the ACTG~175
composite.

\subsubsection*{ACTG~175 (RCT) side}

ACTG~175 releases each subject's endpoint pre-assembled as a pair
$(\texttt{days}, \texttt{cens})$, where $\texttt{cens} \in \{0, 1\}$ is
the composite indicator for the first occurrence of any of (i)~confirmed
$\geq 50\%$ decline in CD4 count from baseline, (ii)~AIDS-defining
diagnosis, or (iii)~death from any cause; and \texttt{days} is the time
from randomization to that event or to administrative censoring, with a
maximum follow-up of approximately $1{,}231$ days. The three components
are not released separately, so the RCT side cannot be respecified as a
subset of the composite. On the analysis scale we work with
$T = \texttt{days}/365.25$ in years and $\delta = \texttt{cens}$.

A separate CD4-50\% reconstruction is also computed from the trial
release (\texttt{cd\_decl\_50}, \texttt{cd\_decl\_days}) by comparing the
20- and 96-week CD4 measurements to half of baseline and approximating
the event time by the nominal window midpoint ($140$ days for the
20-week visit, $672$ days for the 96-week visit). This is reported as a
diagnostic only: it is unconfirmed, ignores AIDS and death, and is never
used as the analysis endpoint.

\subsubsection*{MACS (RWD) side: per-subject components}

The MACS base frame, constructed once per analysis, carries every
per-subject component as either a calendar year or a binary flag,
independent of which variant is later assembled. The components and
their sources are summarized in Table~\ref{tab:macs-components}.

\begin{table}[ht]
  \centering
  \small
  \begin{tabular}{@{}lll@{}}
    \toprule
    Field & Source & Meaning \\
    \midrule
    \texttt{anchor\_year} & $\min(\texttt{AVQY})$ in 1991--95 with
        $\texttt{AVNW} = 2$ &
        time-zero (year resolution) \\
    \texttt{aids\_yr}     & \texttt{outcome.DATE1yy} if
        $\texttt{AIDSCASE} \in \{2, 3\}$ &
        year of CDC AIDS diagnosis \\
    \texttt{cd4\_abs\_yr} & \texttt{outcome.SELFCD4Dyy} if $> 0$ &
        year of first absolute CD4 $< 200$ or $< 14\%$ \\
    \texttt{cd4\_rel\_yr} & first $\texttt{LDATY}$ with
        $\texttt{LEU3N} \leq 0.5 \cdot \texttt{cd4\_base}$ &
        year of first relative $\geq 50\%$ CD4 decline \\
    \texttt{cd4\_base}    & first 1991--95 \texttt{LEU3N} per subject &
        baseline CD4 for the relative rule \\
    \texttt{death\_any}   & $\texttt{DEATH} \in \{1, 2, 3, 4\}$ &
        all-cause death indicator (date redacted) \\
    \texttt{death\_aids}  & $\texttt{DEATH} \in \{1, 2\}$ &
        AIDS-related death indicator (date redacted) \\
    \texttt{last\_year}   & $\max(\texttt{LDATY})$ in \texttt{lab\_rslt} &
        last contact year; death-time fallback / censoring \\
    \bottomrule
  \end{tabular}
  \caption{MACS per-subject component fields. Years are calendar years
    extracted from the public release; flags are derived from the
    \texttt{outcome} form. \texttt{LEU3N} is the \texttt{lab\_rslt} CD4
    count, \texttt{LDATY} the lab visit year, \texttt{AVQY} the
    drug-form visit year, and \texttt{AVNW} the ``currently taking''
    indicator on a drug record.}
  \label{tab:macs-components}
\end{table}

Two construction choices are non-trivial.

\paragraph{CD4 reconstruction.} The ACTG~175 $\geq 50\%$-decline rule
does not match the \texttt{SELFCD4D} field in MACS, which encodes an
absolute sub-threshold crossing ($< 200$~cells/$\mu$L or $< 14\%$). We
therefore rebuild the relative rule per subject from the
\texttt{lab\_rslt} time series: take the earliest 1991--95 \texttt{LEU3N}
as $\texttt{cd4\_base}$, then take \texttt{cd4\_rel\_yr} as the year of
the first subsequent lab with $\texttt{LEU3N} \leq 0.5 \cdot
\texttt{cd4\_base}$. This reconstruction is a one-shot threshold
crossing rather than a confirmed decline, which is a known asymmetry
against the RCT-side endpoint and is sensitive to non-uniform
CD4-measurement frequency between subjects.

\paragraph{Missing death dates.} The field \texttt{DTHDATEyy} is redacted
from the public release. We therefore retain death indicators
(\texttt{death\_any}, \texttt{death\_aids}) but no death dates. For
subjects flagged as deceased without a preceding AIDS or CD4 year, the
event-time fallback is the subject's \texttt{last\_year}, i.e.\ the year
of their last laboratory visit. This is a downward bias on death-event
times that propagates into the upper tail of the survival distribution.

A prior-AIDS exclusion is applied as a final step: subjects with
$\texttt{aids\_yr} < \texttt{anchor\_year}$ are dropped so that the
cohort is AIDS-free at time-zero, matching the trial's entry criterion.

\subsubsection*{Assembly}

Given the base frame and a boolean specification of which components
feed the composite, the assembler $\texttt{combine\_macs}(\cdot)$
computes
\begin{align*}
  \texttt{ev\_yr}
    &= \min \bigl\{
        \texttt{cd4\_abs\_yr},\;
        \texttt{cd4\_rel\_yr},\;
        \texttt{aids\_yr},\;
        \texttt{last\_year} \cdot \mathbf{1}\{\text{death active}\}
      \bigr\}_{\text{activated}} , \\
  \texttt{end\_year}
    &= \begin{cases}
        \texttt{ev\_yr}     & \text{if any activated component fires,} \\
        \texttt{last\_year} & \text{otherwise (right-censored at last contact),}
      \end{cases} \\
  T  &= \max \bigl( \texttt{end\_year} - \texttt{anchor\_year},\, 0.5
       \bigr), \\
  \delta &= \mathbf{1}\{\text{any activated component fires}\}.
\end{align*}
Non-fired components contribute $+\infty$ to the rowwise minimum and so
are ignored. Death contributes its fallback year (\texttt{last\_year})
to the minimum only when the chosen death flag is active and the
subject is flagged dead; otherwise the death component does not enter
at all. The half-year floor on $T$ accommodates year-resolution event
times when event and anchor coincide, where the elapsed time would
otherwise round to exactly zero and trigger a $\log(0)$ on the
log-survival scale used downstream.

\subsubsection*{Variant menu}

Five composite specifications are tabulated to map the activation
choices onto recognizable named endpoints
(Table~\ref{tab:variants}).

\begin{table}[ht]
  \centering
  \small
  \begin{tabular}{@{}lll@{}}
    \toprule
    Variant & Name & Components combined (event $=$ earliest activated year) \\
    \midrule
    v1 & \texttt{EFS-current}   &
        absolute CD4 $+$ AIDS $+$ AIDS-related death \\
    v2 & \texttt{EFS-death-fix} &
        absolute CD4 $+$ AIDS $+$ all-cause death \\
    v3 & \texttt{EFS-cd4-fix}   &
        relative CD4 $+$ AIDS $+$ all-cause death \\
    v4 & \texttt{PFS}           &
        AIDS $+$ all-cause death \\
    v5 & \texttt{OS}            &
        all-cause death \\
    \bottomrule
  \end{tabular}
  \caption{MACS endpoint variants. v1 is the original specification; v2
    corrects a death-code restriction that silently dropped roughly
    a quarter of recorded deaths by keeping only AIDS-related death
    codes; v3 swaps the absolute CD4 component for the
    confirmed-decline analogue and is the closest match to the
    ACTG~175 composite. v4 and v5 are progressively narrower
    endpoints used for sensitivity.}
  \label{tab:variants}
\end{table}

\subsubsection*{Which variant the diagnostic figure uses}

The Kaplan--Meier diagnostic figure does not call the variant menu
directly: it reimplements variant v5 (\texttt{OS}) inline. The event is
$\texttt{DEATH} \in \{1, 2, 3, 4\}$, the event time is
\texttt{last\_year} (death-date fallback, since no death date is
recorded), and the censoring time for survivors is also
\texttt{last\_year}. Elapsed time is
$T = \max(\texttt{last\_year} - \texttt{anchor\_year},\, 0.5)$. The
figure labels this curve ``EFS'' for visual comparison with the RCT
side, but mechanically it is overall survival --- a strict subset of
the RCT composite, chosen because death is the only event MACS
captures reliably enough to give a clean survival signal at year
resolution.

\subsubsection*{Time computation, side by side}

Table~\ref{tab:time} summarizes the time fields on the two sides.

\begin{table}[ht]
  \centering
  \small
  \begin{tabular}{@{}lll@{}}
    \toprule
    Quantity & ACTG~175 & MACS \\
    \midrule
    Anchor       & date of randomization (day) &
                   $\min(\texttt{AVQY})$ in 1991--95 with
                   $\texttt{AVNW} = 2$ (year) \\
    Event time   & \texttt{days} since randomization &
                   rowwise $\min$ over activated component years \\
    Censoring    & published in \texttt{days} &
                   \texttt{last\_year} (year of last lab visit) \\
    Analysis $T$ & $\texttt{days} / 365.25$ &
                   $\max(\texttt{end\_year} - \texttt{anchor\_year},\,
                   0.5)$ \\
    Resolution   & day & year (with $0.5$-year floor) \\
    \bottomrule
  \end{tabular}
  \caption{Time-axis construction on the two sides.}
  \label{tab:time}
\end{table}

The $0.5$-year floor on the MACS side biases short-fuse events upward
by up to one year. It is the most visible footprint of the
year-resolution public release on the analysis; its impact concentrates
on subjects whose event year equals their anchor year and is largest
for cohorts enriched in fast events.

\subsubsection*{Cohort definition versus endpoint}

The cohort definition (which subjects enter the at-risk set) is
logically distinct from the endpoint construction (how each subject's
event and time are defined). Changing the cohort --- for example,
restricting to incident users of the contrast regimens within 1991--95
by excluding any subject with a pre-1991 record of
$\texttt{AVNW} = 2$ on a regimen in the contrast code set --- leaves
the endpoint mechanics unchanged. The same event field on a different
at-risk set yields different event counts, different follow-up
distributions, and therefore different survival curves, without any
change to the assembly logic above.

\end{document}
